% precinct-and-band-mapping-examples.tex — Precinct 与 Band 映射算例
% Source: examples/precinct-and-band-mapping-examples.md

\chapter{Precinct 与 Band 映射算例}\label{ex:precinct-band}

对应章节：第 \ref{ch:bands-precincts} 章（子带、Precinct、Packet 与渐进）。

\begin{examplebox}[title={例 1：4:4:4 图像的 IDS 构建}]

\textbf{输入条件}：$16 \times 8$ 图像，3 分量 RGB，DWT 非对称分解。

\medskip
\textbf{已知参数}：

\begin{tabular}{@{}lc@{}}
\toprule
参数 & 值 \\
\midrule
图像尺寸 & $16 \times 8$ \\
分量 & 3（RGB），$s_x = s_y = 1$ \\
$NL_x$ & 2 \\
$NL_y$ & 1 \\
$S_d$ & 0 \\
$C_w$ & 0 \\
\bottomrule
\end{tabular}

\medskip
\textbf{Step 1}：计算 band 数

每分量 band 数：$2 \cdot NL_y + NL_x + 1 = 2 \times 1 + 2 + 1 = 5$。总 band 数：$3 \times 5 = 15$。

\textbf{Step 2}：计算子带尺寸

对于 $NL_x = 2, NL_y = 1$，每个分量的 2D 分解：

\begin{lstlisting}[numbers=none]
Phase 1 (垂直+水平, d=0):
  LL1:  8 x 4
  HL1:  8 x 4
  LH1:  8 x 4
  HH1:  8 x 4

Phase 2 (仅水平, d=1, 对 LL1):
  LL2:  4 x 4
  HL2:  4 x 4
\end{lstlisting}

共 5 个 band/分量：LL2（Phase 2 低频）, HL2（Phase 2 高频）, HL1, LH1, HH1（Phase 1 的三个高频）。注意 Phase 1 产生的 LL1 在 Phase 2 中被进一步分解为 LL2 和 HL2，因此 LL1 本身不作为独立 band 存在。

\textbf{Step 3}：Precinct 划分

$P_y = 2^{NL_y} = 2$（LL 尺度行数）。每 band 的 precinct 高度：

\begin{itemize}
  \item LL2, HL2, HL1：$d_y = 0$，precinct 高度 = $2 / 2^0 = 2$ 行
  \item LH1, HH1：$d_y = 1$，precinct 高度 = $2 / 2^1 = 1$ 行
\end{itemize}

$n_{py} = \lceil 4 / 2 \rceil = 2$ 行 precinct。

\textbf{Step 4}：列划分

$C_w = 0$ 时，$\mathit{cs} = W = 16$。$n_{px} = \lceil 16 / 16 \rceil = 1$ 列（整幅图像为一列）。

\textbf{Step 5}：Packet 顺序（\texttt{pi[]}）

每个 precinct 包含所有 band 行，按 subpacket 组织：

\begin{lstlisting}[numbers=none]
pi[0]:  band=LL2, y=0, subpkt=0
pi[1]:  band=HL2, y=0, subpkt=0
pi[2]:  band=HL1, y=0, subpkt=0
pi[3]:  band=LH1, y=0, subpkt=0
pi[4]:  band=HH1, y=0, subpkt=0
pi[5]:  band=LL2, y=1, subpkt=1
pi[6]:  band=HL2, y=1, subpkt=1
...
(重复 3 个分量)
\end{lstlisting}

\medskip
\textbf{最终结果}：每分量 5 个 band，共 15 个 band。每分量 2 个 precinct 行。LL2/HL2/HL1 各 2 行，LH1/HH1 各 1 行（$d_y=1$），每分量 $3\times2 + 2\times1 = 8$ 个 position。
\end{examplebox}

\begin{examplebox}[title={例 2：4:2:0 图像的 Band 映射}]

\textbf{输入条件}：$16 \times 8$ 图像，3 分量 YUV 4:2:0，DWT 一层分解。

\medskip
\textbf{已知参数}：

\begin{tabular}{@{}lc@{}}
\toprule
参数 & 值 \\
\midrule
图像尺寸 & $16 \times 8$ \\
Y 分量 & $16 \times 8$（$s_x=1, s_y=1$） \\
Cb/Cr 分量 & $8 \times 4$（$s_x=2, s_y=2$） \\
$NL_x$ & 1 \\
$NL_y$ & 1 \\
\bottomrule
\end{tabular}

\medskip
\textbf{分量尺寸}：

\begin{lstlisting}[numbers=none]
comp 0 (Y):  width=16, height=8
comp 1 (Cb): width=8,  height=4
comp 2 (Cr): width=8,  height=4
\end{lstlisting}

\textbf{各分量 LL 子带}：

\begin{lstlisting}[numbers=none]
comp 0 LL: 8 x 4
comp 1 LL: 4 x 2
comp 2 LL: 4 x 2
\end{lstlisting}

\textbf{Precinct 尺寸}：$P_h = 2^{NL_y} = 2$（LL 尺度行数）。

comp 0 的 $n_{py} = \lceil 4 / 2 \rceil = 2$。

comp 1,2 的 $n_{py} = \lceil 2 / 2 \rceil = 1$。

注意：不同分量可能有不同数量的 precinct 行，需要正确处理边界。

\medskip
\textbf{最终结果}：3 分量共 12 个 band（LL/HL/LH/HH $\times$ 3 分量）。comp 0 有 2 个 precinct 行，comp 1,2 各 1 个。
\end{examplebox}

\begin{examplebox}[title={例 3：Precinct 系数提取}]

\textbf{输入条件}：1 个分量，$4 \times 4$ 图像，展示 $F_q$ 定点转换和 GCLI 计算。

\medskip
\textbf{已知参数}：

\begin{tabular}{@{}lc@{}}
\toprule
参数 & 值 \\
\midrule
分量数 & 1 \\
图像尺寸 & $4 \times 4$ \\
$NL_x, NL_y$ & 1, 1 \\
$F_q$ & 4 \\
子带尺寸 & 各 $2 \times 2$ \\
\bottomrule
\end{tabular}

\medskip
\textbf{Step 1}：DWT 系数（以 LL 子带为例）

\begin{lstlisting}[numbers=none]
LL 子带 (2x2):
[240, 180]
[200, 220]
\end{lstlisting}

\textbf{Step 2}：Sign-Magnitude 转换

$F_q = 4$，$F_{q,r} = (1 \ll 4) \gg 1 = 8$

\begin{lstlisting}[numbers=none]
240 -> (240 + 8) >> 4 = 15
180 -> (180 + 8) >> 4 = 11
200 -> (200 + 8) >> 4 = 13
220 -> (220 + 8) >> 4 = 14
\end{lstlisting}

Sign-Magnitude 结果：$[15, 11, 13, 14]$

\textbf{Step 3}：GCLI 计算

\begin{lstlisting}[numbers=none]
Group 0 (4 个系数): or_all = 15 | 11 | 13 | 14 = 0b1111 = 15
GCLI = BSR(15) + 1 = 3 + 1 = 4
\end{lstlisting}

含义：这 4 个系数都需要 4 bit 表示。

\textbf{Step 4}：逆过程（\texttt{precinct\_to\_image()}）

\begin{lstlisting}[language=C,numbers=none]
// sign-magnitude -> signed, 左移 Fq
if (sm & SIGN_BIT_MASK)
    val = -(sm & ~SIGN_BIT_MASK) << Fq
else
    val = sm << Fq
\end{lstlisting}

$15 \ll 4 = 240$，完美恢复原始值。

\medskip
\textbf{最终结果}：$F_q=4$ 时 LL 子带系数 $[240, 180, 200, 220]$ 转换为 sign-magnitude $[15, 11, 13, 14]$，GCLI=4。逆转换完美恢复原始值。
\end{examplebox}
